\chapter{Conceito do projeto}
\label{chap:fundteor}
%--------- NEW SECTION ----------------------
Os robôs móveis têm a capacidade de se moverem sem a assistência de um operador humano. Os mesmos podem ser classificados,quanto ao sistema de locomoção, como terrestres, aquáticos e aéreos. Os terrestres são subdivididos em robôs que possuem rodas, pernas (bípedes) ou esteiras (ref:Review ArticleA review of mobile robots:). Cada um desses métodos possuem características especifícas quanto ao movimento a ser realizado. Os bípedes, por exemplo, simulam um caminhar antropomórfico, semelhante aos humanos.

Conforme destado no Capítulo \ref{chap:intro}

O desenvolvimento deste projeto consiste em produzir um robô que possa caminhar sobre duas pernas. Além disso, o walker deve se locomover de forma autonôma a fim de realizar uma dada missão.

Neste capítulo serão abordados os requisitos do cliente, os requisistos técnicos, a missão do robô e a pesquisa por similares. 



%conferir se precisa de requisitos do cliente
\section{Requisitos do cliente}
 O cliente definiu certos requisitos quanto à operação e  às características do robô:
 \begin{itemize}
    \item Operar em uma área de 2x1,5m;
    \item Possuir uma altura de aproximadamente 30 cm;
    \item Ser capaz de operar por, no mínimo 20 minutos;
    \item Ser capaz de desviar de obstáculos;
 \end{itemize}

 \section{Requisitos técnicos}


   \lipsum[2-4]

 \section{Missão}
 \lipsum
 %desenvolver mais
 Além disso, o Walker deve realizar um desafio, que consiste em navegar de forma autônoma, se localizar por meio de tags e encontrar um determinado objeto.


 \section{Pesquisa por similares}

\lipsum[1]

 \subsection{Carro Voador}
 \lipsum[2-4]
%----------------------------------------------------------

%--------- NEW SECTION ----------------------


%---------------picture------------------------------------
% \begin{figure}
%     \centering
%     \subfigure[Figure A]{\label{fig:a}\includegraphics[width=60mm]{./lq}}
%     \subfigure[Figure B]{\label{fig:b}\includegraphics[width=60mm]{./lq}}
%     \subfigure[Figure C]{\label{fig:c}\includegraphics[width=\textwidth]{./lq}}
%     \caption{Three simple graphs}
%     \label{fig:three graphs}
% \end{figure}
%----------------------------------------------------------

% \begin{figure}
%     \centering
%     \begin{subfigure}[b]{0.3\textwidth}
%         \centering
%         \includegraphics[width=\textwidth]{./lq}
%         \caption{$y=x$}
%         \label{fig:y equals x}
%     \end{subfigure}
%     \hfill
%     \begin{subfigure}[b]{0.3\textwidth}
%         \centering
%         \includegraphics[width=\textwidth]{./lq}
%         \caption{$y=3sinx$}
%         \label{fig:three sin x}
%     \end{subfigure}
%     \hfill
%     \begin{subfigure}[b]{0.3\textwidth}
%         \centering
%         \includegraphics[width=\textwidth]{./lq}
%         \caption{$y=5/x$}
%         \label{fig:five over x}
%     \end{subfigure}
%        \caption{Three simple graphs}
%        \label{fig:three graphs}
% \end{figure}


% %--------- NEW SECTION ----------------------
% \section{Assunto 2}
% \label{sec:ass2}
% flkjasdlkfjasdlkfjs

% \begin{table}[h]
%     \begin{subtable}[h]{0.45\textwidth}
%         \centering
%         \begin{tabular}{l | l | l}
%         Day & Max Temp & Min Temp \\
%         \hline \hline
%         Mon & 20 & 13\\
%         Tue & 22 & 14\\
%         Wed & 23 & 12\\
%         Thurs & 25 & 13\\
%         Fri & 18 & 7\\
%         Sat & 15 & 13\\
%         Sun & 20 & 13
%        \end{tabular}
%        \caption{First Week}
%        \label{tab:week1}
%     \end{subtable}
%     \hfill
%     \begin{subtable}[h]{0.45\textwidth}
%         \centering
%         \begin{tabular}{l | l | l}
%         Day & Max Temp & Min Temp \\
%         \hline \hline
%         Mon & 17 & 11\\
%         Tue & 16 & 10\\
%         Wed & 14 & 8\\
%         Thurs & 12 & 5\\
%         Fri & 15 & 7\\
%         Sat & 16 & 12\\
%         Sun & 15 & 9
%         \end{tabular}
%         \caption{Second Week}
%         \label{tab:week2}
%      \end{subtable}
%      \caption{Max and min temps recorded in the first two weeks of July}
%      \label{tab:temps}
% \end{table}